\chapter{Transaction-Cost Analysis}\label{ch:mx:transaction-cost-analysis}

A portfolio manager's order cost 25 basis points against the arrival price; the broker's report says it beat VWAP by three. Both numbers are right, and
only one of them says what the order cost the fund. Transaction-cost analysis (One Quant Book 7, chapter~23) is the discipline of choosing the number
that answers the question asked, of explaining it, and of comparing orders that were not equally hard. This chapter builds a simulated order log in
which every order was worked by two algorithms, so that the true difference between them is known, and then analyses the log a desk would actually
have, in which each order saw only one.

\section{Benchmarks and what each hides}

\begin{definition}[Execution benchmark]\label{def:mx:transaction-cost-analysis:benchmark}
An \emph{execution benchmark}\index{execution benchmark} is the reference price an execution is measured against: the decision or arrival price
(implementation shortfall), the volume-weighted price over the order's interval or the day, the close. The slippage against it is the average
execution price minus the benchmark, signed so that a positive number is a cost.
\end{definition}

The simulated log has 100 days of \texttt{firm.agentmkt} on \texttt{firm.exchsim}, each 18 minutes long, with a volatility regime (the hidden value
jumping 0.1, 0.2 or 0.4 times a second) and an activity level (0.7, 1 or 1.4 times the base flow) drawn per day. Each day carries two \omterm{def:mx:the-almgren-chriss-framework:parent}{parent orders},
a buy or a sell of 2\,000 to 12\,000 shares (2\% to 24\% of the volume over their window), decided 30 seconds before they start, worked over five minutes
in ten 30-second slices, with a limit 20 ticks from the arrival price. Two algorithms (chapter~16's schedules on chapter~17's executor) work them:
\begin{itemize}
\item \emph{patient}: equal slices, each resting at the best quote and repriced, with a \omterm{def:mx:child-order-placement:risk}{clean-up trade} at the slice's end;
\item \emph{urgent}: an Almgren--Chriss schedule front-loaded with $\kappa T=3$ (chapter~14), every slice crossed at once.
\end{itemize}
Every order is worked by both, and the day is also run without it, all with the same exogenous order flow (chapter~11's counterfactual). Prices are
in ticks: a tick is a cent on a USD~100 stock, about a basis point.

\begin{table}[ht]
\centering
\small
\begin{tabular}{@{}lrrrr@{}}
\toprule
algorithm & vs arrival & vs interval VWAP & vs day VWAP & vs close \\
\midrule
patient & 2.42 & $-0.29$ & 1.90 & 2.51 \\
urgent & 3.66 & 1.40 & 3.39 & 4.10 \\
\bottomrule
\end{tabular}
\caption{Average slippage per filled share (ticks) of the 200 simulated orders under each algorithm, against four benchmarks. Data:
\texttt{mx\_tca.study}.}\label{tab:mx:transaction-cost-analysis:benchmarks}
\end{table}

The patient algorithm ``beats'' the interval VWAP by 0.29 ticks while costing 2.42 against the arrival price: the VWAP over its own window includes its
own trades and the price they pushed, so it moves with the order. Against the day's VWAP or the close, both algorithms look better or worse depending
on what the rest of the day did. Only the arrival price, or the decision price before it, is fixed before the order touches the market. Perold (1988)
set the paper portfolio, traded at the decision price, against the real one: the implementation shortfall.

\section{Pre-trade estimates}

\begin{definition}[Pre-trade cost estimate]\label{def:mx:transaction-cost-analysis:pretrade}
A \emph{pre-trade cost estimate}\index{pre-trade cost estimate} is the expected cost of an order computed before it is sent, from its size, the
stock's volume, volatility and spread, and a model of impact; it prices the decision to trade and sets the yardstick for the result.
\end{definition}

The square-root law (One Quant Book 7, chapter~27, and chapter~11) gives the model: a half-spread plus $\eta\,\sigma\sqrt{Q/V}$, with $\sigma$ the price's
volatility over the order's horizon and $Q/V$ its participation. The desk estimates $\sigma$ for each volatility regime from its days without orders
(7.6, 10.0 and 11.5 ticks over five minutes) and fits $\eta$ on the patient algorithm's orders of the first 50 days (with $-0.5$ tick, the half-spread
it earns, as the intercept): $\eta=1.04$ with a standard error of 0.18 (\texttt{firm\_tca.fit\_pretrade}, on Book~7's \texttt{firm.tcost}). On the
next 50 days the model predicts 2.47 ticks a share on average and the orders cost 2.73 (standard error 0.52): right on average, but the correlation
between prediction and outcome across orders is only 0.25 (\cref{fig:mx:transaction-cost-analysis:calibration}). A single order's cost is mostly the
price's own move; a pre-trade model is a statement about averages.

\begin{omfigure}
\begin{tikzpicture}
\begin{axis}[width=8.6cm,height=5.8cm,xlabel={\small predicted cost (ticks a share)},ylabel={\small realised cost (ticks a share)},
  tick label style={font=\footnotesize},grid=major,grid style={gray!25},xmin=0,xmax=5,ymin=-2,ymax=8,table/col sep=comma]
\addplot[omDef,only marks,mark=*,error bars/.cd,y dir=both,y explicit] table[x=pred,y=real,y error=se]{figdata/microstructure/19-transaction-cost-analysis/calibration.csv};
\addplot[black,dashed,domain=0:5] {x};
\end{axis}
\end{tikzpicture}

\omcaption{The pre-trade model out of sample: the patient algorithm's orders of the last 50 days in five groups by predicted cost, their mean realised
cost against arrival and its standard error. Data: \texttt{mx\_tca.study}.}\label{fig:mx:transaction-cost-analysis:calibration}
\end{omfigure}

\section{Post-trade attribution}

\begin{definition}[Slippage attribution]\label{def:mx:transaction-cost-analysis:attribution}
\emph{Slippage attribution}\index{slippage attribution} splits an order's implementation shortfall into parts that add up to it: the delay before the
order reaches the market, the spread paid or earned on each fill, the impact (the price's move caused by the order), the timing (the price's move it
did not cause), the opportunity cost of what was not filled, and fees.
\end{definition}

With a decision price $d$, arrival $a$, fills $q_i$ at $p_i$ when the mid was $m_i$ (just before the fill) and would have been $\tilde m_i$ without the
order, target $X$ and final mid $e$, the shortfall of a buy is
\[
\underbrace{\textstyle\sum q_i(a-d)}_{\text{delay}}+\underbrace{\textstyle\sum q_i(p_i-m_i)}_{\text{spread}}+\underbrace{\textstyle\sum q_i(m_i-\tilde m_i-(a-\tilde a))}_{\text{impact}}+\underbrace{\textstyle\sum q_i(\tilde m_i-\tilde a)}_{\text{timing}}+\underbrace{(X-\textstyle\sum q_i)(e-d)}_{\text{opportunity}}+\text{fees},
\]
which is Perold's delay, execution, opportunity and fees (\texttt{firm\_markout.shortfall}) with the execution split in three. Real TCA does not see
$\tilde m$; it splits impact from timing with a model or with the market's move. The simulator sees it, which is what makes the split below exact.

\omcode{code/firm/tca/firm_tca.py}{44}{60}{Attribution of a parent order's implementation shortfall, per share of its target: the mid before each fill against the fill price (spread), against the mid the same market had without the order (impact), and that mid's own move (timing).}\label{lst:mx:transaction-cost-analysis:attribute}

\begin{omfigure}
\begin{tikzpicture}
\begin{axis}[width=11cm,height=5.6cm,ybar,bar width=8pt,ylabel={\small ticks per share of target},
  xtick={0,1,2,3,4,5,6},xticklabels={delay,spread,impact,timing,{opport.},fees,total},
  tick label style={font=\footnotesize},ymin=-1.5,ymax=4.5,grid=major,grid style={gray!25},enlarge x limits=0.08,
  legend columns=2,legend style={font=\footnotesize,at={(0.5,-0.2)},anchor=north,draw=none,fill=none,
  /tikz/every even column/.append style={column sep=8pt}},table/col sep=comma]
\addplot[fill=omDef!70,draw=omDef] table[x=i,y=patient]{figdata/microstructure/19-transaction-cost-analysis/attribution.csv};
\addplot[fill=omProp!70,draw=omProp] table[x=i,y=urgent]{figdata/microstructure/19-transaction-cost-analysis/attribution.csv};
\legend{patient,urgent}
\end{axis}
\end{tikzpicture}

\omcaption{Attribution of the implementation shortfall of the 200 orders under each algorithm: the parts add up to the total exactly. Data:
\texttt{mx\_tca.study}.}\label{fig:mx:transaction-cost-analysis:attribution}
\end{omfigure}

The two algorithms spend differently (\cref{fig:mx:transaction-cost-analysis:attribution}). The patient one earns 0.65 ticks of spread and pays 2.41 of
impact: resting a large order at the best quote for five minutes holds the price on its side. The urgent one pays 2.43 of spread, which includes walking
the book on its large early slices, and 1.13 of impact. Timing averages near zero (0.59 and $-0.07$), as it should with random sides; opportunity is
small (99.5\% and 98.0\% filled within the limit); fees are a rebate for the patient algorithm ($-0.20$) and a charge for the urgent one (0.29). In
total, 1.90 against 3.79 ticks per share, with standard deviations of 5.98 and 4.60: the urgent algorithm costs more and varies less.

\begin{definition}[Post-trade reversion]\label{def:mx:transaction-cost-analysis:reversion}
\emph{Post-trade reversion}\index{post-trade reversion} is the move of the price after an order's last fill, signed against the order: a price that
gives back the order's push (negative) shows temporary impact; a price that keeps going shows that the order carried information or that it leaked
(chapter~9).
\end{definition}

Two minutes after the last fill the price has given back 3.51 ticks for the patient algorithm and 1.06 for the urgent one (1.37 and 0.07 after 30
seconds). The patient algorithm's last fills are often its \omterm{def:mx:child-order-placement:risk}{clean-up trades}, which push the price at the end of the window; the urgent algorithm's last
slices are small. No simulated order carries information, so no price keeps going; in real logs, a continuation after completion is the leakage signal.

\section{Comparing across orders, brokers and peers}

Which algorithm is cheaper? In the full log the answer is known: working every order with both, the urgent algorithm costs 1.89 ticks a share more,
with a standard error of 0.25 clustered by day. The desk's log is different. Its routing gave the urgent algorithm large orders on volatile days more
often (53.5\% of orders went to it, averaging 7\,039 shares against 4\,126 for the patient one), so the raw difference in its log is 3.36 ticks
(standard error 0.76): hard orders made the urgent algorithm look worse than it is.

\begin{definition}[Peer-universe comparison]\label{def:mx:transaction-cost-analysis:peer}
A \emph{peer-universe comparison}\index{peer-universe comparison} measures an execution's cost against that of comparable orders from other traders,
brokers or funds, after adjusting for how hard each order was (size against volume, volatility, spread, \omterm{def:mx:the-almgren-chriss-framework:urgency}{urgency}), usually as the residual of a cost
model fitted on the universe.
\end{definition}

Adjusting means regressing the cost on the algorithm and on the order's difficulty (the regime's volatility, the participation, the square-root
term and its interaction with the algorithm), with standard errors clustered by day because the two orders of a day share its market (clustered errors are One Quant Book 4,
chapter~16's); \texttt{firm\_tca.adjusted\_difference} returns both the raw and the adjusted coefficient.

\omcode{code/microstructure/19-transaction-cost-analysis/python/mx_tca.py}{148}{152}{The desk's routing in the simulated log: the probability of sending an order to the urgent algorithm rises with its size and on volatile days, which is what makes the raw comparison biased.}\label{lst:mx:transaction-cost-analysis:assigned}

\begin{omfigure}
\begin{tikzpicture}
\begin{axis}[width=9cm,height=4.4cm,xlabel={\small urgent minus patient (ticks a share)},xmin=-0.5,xmax=5.5,ymin=0.4,ymax=3.6,
  ytick={1,2,3},yticklabels={{truth (both algorithms)},{adjusted},{raw}},tick label style={font=\footnotesize},grid=major,grid style={gray!25},
  table/col sep=comma]
\addplot[omDef,only marks,mark=*,mark size=2.4pt,error bars/.cd,x dir=both,x explicit]
  table[x=est,y=y,x error plus expr=\thisrow{hi}-\thisrow{est},x error minus expr=\thisrow{est}-\thisrow{lo}]{figdata/microstructure/19-transaction-cost-analysis/estimates.csv};
\addplot[black,dashed] coordinates {(0,0.4) (0,3.6)};
\end{axis}
\end{tikzpicture}

\omcaption{The cost difference between the two algorithms with 95\% intervals: raw on the desk's log, adjusted for difficulty on the same log, and the
truth from working every order with both. Data: \texttt{mx\_tca.study}.}\label{fig:mx:transaction-cost-analysis:estimates}
\end{omfigure}

The adjusted difference is 2.00 ticks, with a 95\% interval from 0.32 to 3.68 (\cref{fig:mx:transaction-cost-analysis:estimates}): the adjustment
removes most of the bias (3.36 raw against 1.89 true) and the interval covers the truth, but it is wide, because 200 orders with a cost standard
deviation of five or six ticks do not pin down a two-tick difference. Measured against the pre-trade model instead, the patient algorithm's orders
cost 0.77 ticks less than predicted (standard error 0.56) and the urgent one's 1.50 more (0.46): the same conclusion, with the model's own error on top.
Anand, Irvine, Puckett and Venkataraman (2012) found in institutional trade data that trading desks sustain their relative performance from one
period to the next; Frazzini, Israel and Moskowitz (2018), on a large manager's live executions, found real trading costs an order of magnitude
smaller than earlier studies suggested. Both needed difficulty adjustment of this kind before any comparison.

\begin{dated}{2026-09}{Best-execution reports in Europe}\label{dat:mx:transaction-cost-analysis:rts}
MiFID II required execution venues to publish execution-quality reports (``RTS~27'') and investment firms to publish their top five venues each year
(``RTS~28''). Directive (EU) 2021/338 suspended the venues' periodic reports until 28 February 2023. Directive (EU) 2024/790 of 28 February 2024, noting
that the reports were rarely read and did not allow meaningful comparisons, deleted the top-five-venue report and replaced the execution-quality report
with a duty to tell each client where its order was executed; member states had until 29 September 2025 to transpose it. In the United Kingdom the
FCA removed both reports from 1 December 2021 (policy statement PS21/20).
\end{dated}

\section{Closing the loop}

TCA pays for itself when it changes what the desk does. The loop here fits the cost on the algorithm and its interaction with difficulty on the first
50 days of the desk's log, and routes each order of the last 50 days to the algorithm with the lower prediction. The model sends every order to the
patient algorithm; scored with the truth, that costs 2.17 ticks a share on those days, against 3.24 for the desk's routing and 3.88 for sending
everything to the urgent algorithm. The urgent algorithm's lower dispersion is the reason a desk might still choose it: a manager who penalises the
variance of cost enough (exercise~6) pays the 1.9 ticks gladly. The comparison of algorithms by experiment rather than by regression, with orders
assigned at random, is chapter~20's algo wheel; the refinements of A/B testing (One Quant Book 7, chapter~21, including variance reduction by CUPED)
apply directly.

\section{Tutorial: analysing an order log}\label{tut:mx:transaction-cost-analysis:tut}

\begin{tutorial}
\textbf{Goal.} Generate a parent-order log in the simulated market, run TCA against every benchmark, attribute the shortfall, fit and test a pre-trade
model, and compare two algorithms with and without difficulty adjustment against the truth.
\textbf{End state:} \cref{tab:mx:transaction-cost-analysis:benchmarks}, \cref{fig:mx:transaction-cost-analysis:calibration,fig:mx:transaction-cost-analysis:attribution,fig:mx:transaction-cost-analysis:estimates}
and the numbers of sections 4 and 5.
\begin{tutsteps}
\item \textbf{Log.} \texttt{mx\_tca.conditions(d)}, \texttt{orders(d)}, \texttt{day\_runs(d)} (the day without orders, patient, urgent),
\texttt{log()}.
\item \textbf{Measure.} \texttt{firm\_tca.attribute}, \texttt{benchmarks}, \texttt{reversion}.
\item \textbf{Pre-trade.} \texttt{fit\_pretrade}, \texttt{pretrade}.
\item \textbf{Compare.} \texttt{assigned(r)}, \texttt{adjusted\_difference}, \texttt{peer\_compare}; \texttt{study()}; draw with \texttt{fig\_tca.py}.
\end{tutsteps}
\textbf{What to change next.} Give some orders information (a price drift in their direction) and watch reversion turn into continuation; route at
random and see the adjustment become unnecessary; add a third algorithm.
\end{tutorial}

\section{Build: transaction-cost analysis}\label{bld:mx:transaction-cost-analysis:tca}

\begin{build}
\bfield{Purpose} The measurement layer behind the algo wheel of chapter~20, the \omterm{def:mx:benchmark-algorithms:algo}{execution algorithm} of chapter~28 and the desk's reports.
\bfield{Interface} \texttt{attribute(side, target, decision, arrival, fills, mid, cf\_mid, end, t\_arrival)}, \texttt{benchmarks(side, fills, arrival,
window, day, close)}, \texttt{reversion(side, t\_end, mid, horizons)}, \texttt{pretrade(half\_spread, eta, sigma, participation)},
\texttt{fit\_pretrade}, \texttt{cluster\_ols(y, X, clusters)}, \texttt{adjusted\_difference(cost, treat, controls, clusters)},
\texttt{peer\_compare(cost, predicted, groups, clusters)}.
\bfield{Rules} Costs positive when paid; attribution per share of target (opportunity counts), benchmarks per filled share; the mid just before each
fill; clustered errors by the unit that shares a market (the day).
\bfield{Acceptance tests} \texttt{code/firm/tca/tests/}: the attribution on a hand example, adding up to Book~7's shortfall; benchmarks and reversion;
the clustered adjustment recovers a planted effect that the raw difference overstates; the peer comparison.
\bfield{Stretch} Attribution without a counterfactual (impact from the pre-trade model); peer universes by size and liquidity bucket.
\end{build}

\begin{omsources}
\item A.~F.~Perold, ``The implementation shortfall: paper versus reality'', \emph{Journal of Portfolio Management} 14(3), 1988.
\item A.~Anand, P.~Irvine, A.~Puckett and K.~Venkataraman, ``Performance of institutional trading desks: an analysis of persistence in trading
costs'', \emph{Review of Financial Studies} 25(2), 2012.
\item R.~Kissell, \emph{The Science of Algorithmic Trading and Portfolio Management}, Academic Press, 2014.
\item A.~Frazzini, R.~Israel and T.~J.~Moskowitz, ``Trading costs'', working paper, SSRN 3229719, 2018.
\item Directive (EU) 2021/338 and Directive (EU) 2024/790 amending Directive 2014/65/EU (MiFID II); FCA, PS21/20, 2021.
\end{omsources}

\section{Exercises}

\begin{exercise}[$\star$]\label{exo:mx:transaction-cost-analysis:1}
A buy of 1\,000 shares is decided at 100.00 and arrives at 100.02; 800 shares fill at an average of 100.05, the price ends at 100.10, and fees are 0.3
cents a share. Compute the delay, execution, opportunity and fee costs and the shortfall per share of target.
\end{exercise}

\begin{exercise}[$\star$]\label{exo:mx:transaction-cost-analysis:2}
The same order's interval VWAP was 100.06. What does the VWAP report say, and what does the arrival report say?
\end{exercise}

\begin{exercise}[$\star$]\label{exo:mx:transaction-cost-analysis:3}
With the fitted model ($-0.5$ tick of spread, $\eta=1.04$), what is the pre-trade estimate for an order at 10\% participation in the middle regime
($\sigma=10.0$ ticks)?
\end{exercise}

\begin{exercise}[$\star\star$]\label{exo:mx:transaction-cost-analysis:4}
Why is the raw difference between the algorithms in the desk's log larger than the truth?
\end{exercise}

\begin{exercise}[$\star\star$]\label{exo:mx:transaction-cost-analysis:5}
Why is the patient algorithm's reversion larger than the urgent one's, when it trades more gently?
\end{exercise}

\begin{exercise}[$\star\star$]\label{exo:mx:transaction-cost-analysis:6}
A manager minimises expected cost plus $\lambda$ times its variance (per share, in ticks). Above what $\lambda$ does the urgent algorithm's lower
standard deviation justify its higher mean?
\end{exercise}

\begin{exercise}[$\star\star\star$]\label{exo:mx:transaction-cost-analysis:7}
\emph{Coding.} Adjust the comparison for participation only. What difference and interval do you get, and why?
\end{exercise}

\begin{exercise}[$\star\star\star$]\label{exo:mx:transaction-cost-analysis:8}
\emph{Find the flaw.} ``The patient algorithm beat the VWAP by 0.3 ticks on average, so its orders cost our clients nothing.''
\end{exercise}

\section{Problem: Bad Broker or Hard Day?}

\begin{problem}[{Weekend problem --- bad broker or hard day?}]\label{pb:mx:transaction-cost-analysis:1}
A desk's TCA report says the urgent algorithm costs 3.4 ticks a share more than the patient one. The broker behind it says its orders were harder.

\medskip\noindent\textbf{Part I --- Benchmarks.}
\begin{enumerate}
\item Define the four benchmarks of the chapter.
\item Give each algorithm's slippage against them.
\item Why does the patient algorithm beat its interval VWAP?
\item Which benchmark answers ``what did the order cost the fund'', and why?
\end{enumerate}

\medskip\noindent\textbf{Part II --- Estimates and attribution.}
\begin{enumerate}[resume]
\item State the pre-trade model and its fitted coefficient.
\item How well does it predict out of sample, on average and order by order?
\item Write the attribution and show it adds up to the implementation shortfall.
\item Give the attribution of each algorithm.
\item What does the reversion after the last fill say?
\end{enumerate}

\medskip\noindent\textbf{Part III --- The comparison.}
\begin{enumerate}[resume]
\item Describe the desk's routing and the orders each algorithm received.
\item Give the raw difference and its standard error.
\item Write the adjusted regression and explain the clustering.
\item \emph{State the named result}: the difficulty-adjusted cost difference between the two algorithms with its confidence interval, against the
raw difference.
\item How does it compare with the truth, and why can a real desk not see the truth?
\item What does the comparison against the pre-trade model give?
\end{enumerate}

\medskip\noindent\textbf{Part IV --- Closing the loop.}
\begin{enumerate}[resume]
\item What routing does the fitted model recommend, and what does it save?
\item Why might a manager still choose the urgent algorithm?
\item How would randomised routing change the analysis?
\item Bad broker or hard day: what do you answer?
\item In one sentence: what makes two execution costs comparable?
\end{enumerate}
\end{problem}

\section{Interview questions}

\begin{interviewq}[$\star$ \iqroles{trader}]\label{iq:mx:transaction-cost-analysis:1}
Your order beat VWAP but cost 20 basis points against arrival. Explain both numbers to the portfolio manager.
\end{interviewq}

\begin{interviewq}[$\star\star$ \iqroles{researcher}]\label{iq:mx:transaction-cost-analysis:2}
How would you split an order's slippage into spread, impact and timing without a counterfactual?
\end{interviewq}

\begin{interviewq}[$\star\star$ \iqroles{researcher}]\label{iq:mx:transaction-cost-analysis:3}
Two brokers' average costs differ by 3 basis points. What do you need before concluding that one is better?
\end{interviewq}

\begin{interviewq}[$\star\star$ \iqroles{risk}]\label{iq:mx:transaction-cost-analysis:4}
What does \omterm{def:mx:transaction-cost-analysis:reversion}{post-trade reversion} tell you, and how would you use it to detect \omterm{def:mx:dark-pools-and-blocks:leakage}{information leakage}?
\end{interviewq}

\begin{interviewq}[$\star\star$ \iqroles{developer}]\label{iq:mx:transaction-cost-analysis:5}
Design the data a TCA system must capture for each parent and \omterm{def:mx:the-almgren-chriss-framework:parent}{child order}.
\end{interviewq}

\begin{interviewq}[$\star\star\star$ \iqroles{bank}]\label{iq:mx:transaction-cost-analysis:6}
A client's TCA vendor ranks your algorithms last in its peer universe. How do you examine the claim?
\end{interviewq}
